\chapter{Vitamin K}

\section*{What Is This Test?}
The vitamin K test measures the serum or plasma concentration of vitamin K, a fat-soluble vitamin required for the post-translational gamma-carboxylation of glutamic acid residues on a family of proteins that depend on this modification for biological activity. The vitamin K-dependent proteins include coagulation factors II (prothrombin), VII, IX, and X, and the anticoagulant proteins C and S; in bone, the carboxylated proteins osteocalcin and matrix Gla protein depend on the same reaction. Vitamin K exists in two natural forms: phylloquinone (vitamin K1), found in green leafy vegetables and the predominant dietary form, and menaquinones (vitamin K2), produced by intestinal bacteria and found in fermented foods, egg yolk, and meat. Deficiency impairs gamma-carboxylation and produces a bleeding tendency characterized by a prolonged prothrombin time (PT) that corrects with vitamin K administration. Newborns are at particular risk because the placenta transports vitamin K poorly, the neonatal liver has low stores, and the gut is sterile; prophylactic intramuscular vitamin K at birth has virtually eliminated classic hemorrhagic disease of the newborn. Plasma phylloquinone is the direct measure of vitamin K status, while the functional surrogate is the PT/INR, which is prolonged when carboxylation is impaired.

\section*{Alternative Names}
\begin{itemize}
  \item Phylloquinone (Vitamin K1)
  \item Menaquinone (Vitamin K2)
  \item Serum Vitamin K
  \item Vitamin K Status
  \item Prothrombin Time (functional measure)
\end{itemize}
\section*{Principle}
Plasma phylloquinone (K1) is measured by high-performance liquid chromatography (HPLC) with fluorescence detection or by liquid chromatography-tandem mass spectrometry (LC-MS/MS). Because vitamin K concentrations in blood are low (nanogram per milliliter range), the assay requires lipid extraction: plasma proteins are precipitated with ethanol, and vitamin K is extracted into hexane or heptane. The extract is evaporated and reconstituted, then separated on a reverse-phase C18 or C30 column. Post-column reduction with a zinc or platinum reactor converts vitamin K to its fluorescent hydroquinone form for detection at excitation/emission wavelengths of approximately 243/418 nm. Internal standards (deuterated or structurally related compounds) correct for recovery. Menaquinones (K2) are quantified by the same method with appropriate standards. Because the functional status of vitamin K-dependent proteins is assessed by the prothrombin time/INR (clotting assay) and by measurement of PIVKA-II (protein induced by vitamin K absence), vitamin K status is typically evaluated with a combination of the direct concentration and these functional tests.

\section*{History}
The coagulation defect of vitamin K deficiency was described in 1929 by Henrik Dam, who observed a hemorrhagic disease in chicks fed a fat-free diet, and he named the responsible fat-soluble factor ``vitamin K'' (for Koagulation). Dam received the 1949 Nobel Prize in Physiology or Medicine, shared with Edward Doisy, who isolated and synthesized phylloquinone in 1939. The vitamin K cycle and the discovery that dicoumarol (the forerunner of warfarin) inhibits vitamin K epoxide reductase were described by Karl Link in the 1940s. The recognition that vitamin K-dependent proteins require gamma-carboxylation was established by Johan Stenflo in 1974. The practice of routine intramuscular vitamin K prophylaxis for newborns, introduced in the 1960s and reaffirmed by the American Academy of Pediatrics in 2003, dramatically reduced the incidence of hemorrhagic disease of the newborn. Sensitive plasma assays for phylloquinone were developed in the 1970s--1980s and remain the reference for assessing vitamin K status.

\section*{Why This Test Is Ordered}
\begin{itemize}
  \item Evaluation of unexplained bleeding or bruising, especially in patients with malabsorption syndromes, cholestatic liver disease, chronic diarrhea, or prolonged broad-spectrum antibiotic use
  \item Investigation of a prolonged prothrombin time (PT/INR) to determine whether it reflects vitamin K deficiency (corrects within 24 hours of vitamin K administration) or liver disease, warfarin effect, or other coagulation factor deficiency
  \item Assessment of fat-soluble vitamin status in malabsorption syndromes: celiac disease, Crohn disease, short bowel syndrome, cystic fibrosis, chronic pancreatitis, and biliary obstruction
  \item Evaluation of patients on long-term parenteral nutrition or with chronic alcoholism and malnutrition
  \item Newborn evaluation: hemorrhagic disease of the newborn is suspected when bleeding occurs in a neonate who did not receive vitamin K prophylaxis; later-onset disease occurs in breastfed infants with cholestasis or malabsorption
  \item Assessment of patients on chronic antibiotics that deplete vitamin K-producing gut flora
  \item Monitoring of nutritional status in patients on warfarin or in those with suspected vitamin K excess interfering with anticoagulation
  \item Research and selected clinical evaluation of vitamin K-dependent protein status in osteoporosis and vascular calcification
\end{itemize}

\section*{Primary vs Secondary Test}
Plasma phylloquinone is a secondary test. The first-line laboratory evaluation of suspected vitamin K deficiency is the prothrombin time (PT/INR), which is prolonged in deficiency and corrects after vitamin K administration. Plasma vitamin K measurement directly confirms the deficiency state and is used when the diagnosis is uncertain, when functional tests are ambiguous, or in nutritional research.

\section*{How This Test Is Performed}
A blood sample is collected by standard venipuncture into an EDTA (lavender-top) or sodium citrate (light blue-top) tube. The specimen should be protected from light because vitamin K is photolabile, and it must be processed promptly. Plasma is separated and analyzed by HPLC or LC-MS/MS. The PT/INR, when ordered for the same purpose, is measured on citrated plasma by an automated coagulation analyzer. Results for plasma phylloquinone are typically available within 3--7 days (frequently a send-out test); the PT/INR is available within hours.

\section*{What Preparation Is Needed}
Fasting for 8--12 hours is preferred because dietary vitamin K transiently raises plasma levels. The specimen must be protected from light. Patients should disclose all supplements (including multivitamins, which contain phylloquinone) and any warfarin or anticoagulant therapy, because recent vitamin K intake directly raises the measured level and warfarin does not change the plasma concentration but dramatically affects the functional PT. Because vitamin K is fat-soluble, absorption is impaired in fat malabsorption, and the level should be interpreted in that context.

\section*{Contraindications and Precautions}
\begin{itemize}
  \item There are no absolute contraindications to standard venipuncture. Important interpretive cautions:
  \item (1) Plasma phylloquinone reflects recent intake and is not a direct measure of hepatic stores; a normal level does not exclude tissue-level deficiency in the presence of impaired gamma-carboxylation.
  \item (2) Levels fluctuate with diet, so a single fasting sample is the standard, and low-fat diets lower measured values.
  \item (3) Cholestasis and malabsorption lower vitamin K despite normal intake.
  \item (4) The PT/INR is the functional test of choice: it is prolonged in deficiency and corrects within 24 hours of parenteral vitamin K.
  \item (5) PIVKA-II (des-gamma-carboxy prothrombin) is a sensitive marker of subclinical deficiency and of hepatocellular carcinoma.
  \item (6) Specimens exposed to light or left unprocessed degrade rapidly.
  \item (7) In newborns, plasma phylloquinone is low even in health, so reference ranges must be age-specific; the diagnosis of hemorrhagic disease of the newborn is primarily clinical and based on the PT response to vitamin K.
\end{itemize}
\section*{Risks}
Minimal risk. Slight pain or bruising at the venipuncture site. Rarely: hematoma, infection, or vasovagal syncope.

\section*{What Happens After the Test}
Results are available within hours for the PT/INR and within 3--7 days for plasma phylloquinone. If vitamin K deficiency is confirmed or strongly suspected, treatment is oral or parenteral (intravenous or subcutaneous) vitamin K; in life-threatening bleeding, fresh frozen plasma is given while vitamin K takes effect. The PT/INR is repeated within 6--24 hours to document correction. Newborns who did not receive prophylaxis are given intramuscular vitamin K immediately. The underlying cause (malabsorption, liver disease, cholestasis, malnutrition) is then evaluated and treated, and in patients with ongoing risk, maintenance vitamin K supplementation and serial monitoring are arranged. In patients on warfarin, elevated vitamin K intake is managed by consistent dietary vitamin K intake and dose adjustment.

\section*{Understanding Results}

\subsection*{Below Normal}
\begin{itemize}
  \item \textbf{Plasma phylloquinone below the laboratory reference range:} Reflects inadequate intake, malabsorption (celiac disease, Crohn disease, short bowel syndrome, cholestasis, chronic pancreatitis), prolonged antibiotic use reducing gut menaquinone production, or liver disease with impaired transport.
  \item \textbf{Clinical manifestations:} Easy bruising, bleeding from mucosal surfaces, hematuria, gastrointestinal bleeding, and in severe cases intracranial hemorrhage. The PT/INR is prolonged because factors II, VII, IX, and X are incompletely carboxylated; factor VII has the shortest half-life, so the PT rises earliest.
  \item \textbf{Hemorrhagic disease of the newborn:} Early disease occurs within 24 hours (maternal medications such as anticonvulsants or warfarin), classic disease at days 2--7 (low stores, no prophylaxis), and late disease at weeks 2--12 in breastfed infants with cholestasis or malabsorption; intracranial bleeding is a feared complication.
  \item \textbf{Subclinical deficiency:} Elevated PIVKA-II with normal PT/INR indicates partial carboxylation impairment and is associated with reduced osteocalcin carboxylation and, in cohort studies, with increased fracture risk and vascular calcification.
\end{itemize}

\subsection*{Above Normal}
\begin{itemize}
  \item \textbf{Elevated plasma phylloquinone:} Reflects high dietary intake or supplementation; there is no recognized toxicity syndrome for phylloquinone (K1) because it is not stored in toxic amounts, though high doses can temporarily interfere with warfarin anticoagulation.
  \item \textbf{Menadione (synthetic vitamin K3) is not used therapeutically because of toxicity (hemolysis, especially in glucose-6-phosphate dehydrogenase deficiency); its measurement is not part of standard vitamin K panels.}
  \item \textbf{In patients on warfarin, a high plasma vitamin K level from diet or supplements can cause subtherapeutic INRs and increased thrombotic risk; dietary vitamin K should be consistent.}
\end{itemize}

\section*{Normal Results}
\begin{itemize}
  \item \textbf{Plasma phylloquinone (fasting, adults):} 0.1--2.2 ng/mL (0.22--4.9 nmol/L), varying by laboratory
  \item \textbf{PT/INR (functional status):} INR 0.8--1.2 in healthy adults not on anticoagulants; prolonged INRs in deficiency correct with vitamin K
  \item \textbf{PIVKA-II:} normally undetectable or below the laboratory cutoff; elevation indicates impaired gamma-carboxylation
  \item \textbf{Newborn plasma phylloquinone:} typically low (below adult reference values) even in health, rising after the prophylactic dose and with dietary milk intake
  \item \textbf{Daily requirement:} 120 mcg/day for adult men, 90 mcg/day for adult women; newborns receive 0.5--1.0 mg intramuscular prophylaxis at birth
  \item \textbf{Conversion:} 1 ng/mL phylloquinone = 2.22 nmol/L
\end{itemize}

\section*{Follow-up Tests}
\begin{itemize}
  \item \textbf{Prothrombin time/INR:} The key functional test; repeated 6--24 hours after vitamin K administration to document correction
  \item \textbf{PIVKA-II (des-gamma-carboxy prothrombin):} Sensitive marker of subclinical vitamin K deficiency and of hepatocellular carcinoma when markedly elevated
  \item \textbf{Coagulation factor assays (II, VII, IX, X):} When the differential diagnosis includes inherited factor deficiencies rather than vitamin K deficiency
  \item \textbf{Liver function tests and bilirubin:} Distinguish liver disease (which also prolongs the PT but responds poorly to vitamin K) from vitamin K deficiency
  \item \textbf{Other fat-soluble vitamins (A, D, E):} Deficiencies commonly coexist in malabsorption and cholestasis
  \item \textbf{Liver and pancreatic enzyme panels, and imaging if cholestasis or pancreatic insufficiency is suspected:} Identify the cause of fat malabsorption
  \item \textbf{Stool elastase or other pancreatic function tests:} If chronic pancreatitis is suspected as the cause of malabsorption
  \item \textbf{Repeat plasma phylloquinone after replacement:} Confirms repletion in complex nutritional cases
\end{itemize}

\section*{References}
\begin{enumerate}
  \item Shearer MJ, Newman P. Metabolism and cell biology of vitamin K. Thromb Haemost. 2008;100(4):530--547.
  \item Suttie JW. Vitamin K in Health and Disease. CRC Press; 2009.
  \item Shearer MJ, Bolton-Smith C. The U-shaped curve of cardiovascular risk in the elderly: implications of vitamin K. J Intern Med. 2006;260(5):464--472.
  \item American Academy of Pediatrics Committee on Fetus and Newborn. Controversies concerning vitamin K and the newborn. Pediatrics. 2003;112(1 Pt 1):191--192.
  \item Ferland G. Vitamin K and the nervous system: an overview of its actions. Adv Nutr. 2012;3(2):204--212.
  \item Fusaro M, Gallieni M, Rizzo MA, et al. Vitamin K plasma levels determination in human health. Clin Chem Lab Med. 2017;55(6):789--799.
\end{enumerate}

\clearpage
\section*{Regulatory Status and Authorization Date}
This chapter describes a test category, not a specific commercial product. FDA, CDSCO, EU IVDR, and MHRA approval or registration dates therefore cannot be assigned without the assay manufacturer, product name, instrument, and intended use. For laboratory-developed tests, an individual agency approval date may not exist. See the Preface for the interpretation of regulatory status.

\clearpage
